% Esto es una canción típica: empieza con \beginsong{título}[info extra], contiene estrofas
% (enmarcadas en \beginverse y \endverse) y estribillos (\beginchorus y \endchorus).
% Además de la letra, aparecen acordes en la forma \[acorde]. Los acordes se escriben en
% notación inglesa (C-Do, D-Re, E-Mi, F-Fa, G-Sol, A-La, B-Si) y automáticamente se imprimen
% en español (ver archivo estilo/acordes.tex).

% en el bloque de información extra hay varios campos, pero el único que voy a utilizar es
% "by", que indica el autor de la canción.
\beginsong{Bendigamos al Señor}[by={Pelayo Sánchez}]
\beginverse % Empieza una estrofa
Bendi\[E]gamos al Señor,
Dios de \[C#m]toda la creación,
por ha\[A]bernos rega\[B7]lado su a\[E]mor.\[B7]
Su bon\[E]dad y su perdón,
y su \[C#m]gran fidelidad,
por los \[A]siglos de los \[B7]siglos dura\[E]rán.\[E7]
\endverse % Fin de la estrofa
\beginchorus % Empieza un estribillo
El Es\[A]píritu de \[B7]Dios hoy es\[E]tá \[B7]sobre \[C#m]mí,
\[A]Él es quien me ha un\[B7]gido a procla\[E]mar,
\[E7]a procla\[A]mar, la buena \[B7]nueva a los más \[E B7]po\[C#m]bres, \[A]
% La línea anterior es demasiado larga para los ajustes del cancionero, y pasa a la siguiente
% línea automáticamente con un poco de indentación, como si fuera un poema.
% Para ajustar el punto de corte si no cabe todo en una línea, se puede utilizar \brk (dejando espacios a ambos lados)
% en el punto deseado. Por ejemplo "Yo Señor te presento \brk lo que hay en mi interior."
% Por lo general cada línea es un verso.
\ifchorded
% TODO: decidir estilo
(1) la gracia \[E]de su salva\[B7]ción. \rep{2}
(2) la gracia \[B7]de su salva\[E]ción.
\else
la gracia de su salvación. \rep{2}
\fi
% Esto es un condicional: solo se verá lo que hay entre \ifchorded (si tiene acordes) y \fi
% en las versiones del cancionero con acordes, por lo que podemos usarlo para anotaciones
% para los músicos. La estructura {\nolyrics texto con acordes} hace que todo quede en una 
% línea y no en dos líneas como si fuera letra cantable.
\ifchorded
{\nolyrics Puente: \[B7]}
\fi
\endchorus % Fin del estribillo
\beginverse % Empieza otra estrofa: como los acordes son los mismos del ejemplo anterior
% usamos el acento circumflejo para "repetir" el acorde correspondiente. Para más información,
% ver el manual del paquete "songs" en songs.sourceforge.net
Envi^ados con poder
y en el ^nombre de Jesús,
a sa^nar a los en^fermos del do^lor,^
a los ^ciegos dar visión,
a los ^pobres la verdad
y a los ^presos y opri^midos liber^tad.^
\endverse % Fin de la segunda estrofa
\beginverse % Empieza la tercera estrofa
Con la ^fuerza de su amor 
y de ^la resurrección
anun^ciamos llega ^ya la salva^ción, ^
que ni el ^miedo ni el temor,
ni la ^duda o la opresión
borra^rán la paz de ^nuestro cora^zón. ^
\endverse % Fin de la tercera estrofa
\endsong % Fin de la canción (obligatorio)

% Se podrían incluir otras canciones aquí, pero resulta más sencillo separarlas en archivos individuales.
